% Numeric source: paper/shared/evidence/raw/iclr_model_behavior/analysis_v1/results.json
\section{Grounding: Following the Object Across Appearance Changes}
\label{sec:grounding}

Grounding tests should require the model to use both the image and the query.
Our initial target-centered views reveal why this matters: a fixed image-center
point hits every target box in the 30-pair close-view pool. We therefore change the task structure. We place two different
rendered views on one canvas, ask about either category, and swap their positions.
The center is now a neutral gutter between the views. The target location is
determined jointly by the query and the panel arrangement.

\paragraph{Query and position swaps.}
Each board contains two different target categories. Both targets are queried,
and both left/right arrangements are evaluated in each material condition.
Annotations are translated with the corresponding view; boxes are never shown
to the model. Four repaired source pairs form four boards, yielding 16 paired queries per
intervention across two targets and two position permutations. Category queries
ask for the named object. Action queries describe an interaction such as grasping
a cup or turning on a faucet and ask for the corresponding object's location.
The output is a single point on the complete canvas.
Appendix~\ref{app:construction} gives the source pairs and composition details.

\begin{table}[t]
\centering\small
\setlength{\tabcolsep}{3pt}
% Generated by scripts/build_assets.py from corrected results.json
\begin{tabular}{llrrrrr}
\toprule
Model & Intervention & $N$ & Original & Changed & $\Delta$ (pp) & Flips \\
\midrule
Gemma4-E4B & Category query & 16 & 75.0 & 81.2 & +6.2 & 3/16 \\
 & Action query & 16 & 75.0 & 75.0 & +0.0 & 2/16 \\
 & Target window & 16 & 75.0 & 68.8 & -6.2 & 3/16 \\
 & Context window & 16 & 75.0 & 87.5 & +12.5 & 2/16 \\
\midrule
Qwen3-VL-2B & Category query & 16 & 87.5 & 87.5 & +0.0 & 2/16 \\
 & Action query & 16 & 56.2 & 75.0 & +18.8 & 3/16 \\
 & Target window & 16 & 87.5 & 93.8 & +6.2 & 3/16 \\
 & Context window & 16 & 87.5 & 100.0 & +12.5 & 2/16 \\
\midrule
Qwen3-VL-8B & Category query & 16 & 100.0 & 93.8 & -6.2 & 1/16 \\
 & Action query & 16 & 93.8 & 93.8 & +0.0 & 0/16 \\
 & Target window & 16 & 100.0 & 93.8 & -6.2 & 1/16 \\
 & Context window & 16 & 100.0 & 100.0 & +0.0 & 0/16 \\
\bottomrule
\end{tabular}

\caption{Paired two-view interventions. Success rates are percentages over all
queries; invalid outputs count as failures. Original and Changed refer to the
image conditions for the listed intervention. Category and action queries use
full converted views; window interventions exchange only the specified region.
Flips count changes in hit status.}
\label{tab:interventions}
\end{table}

Table~\ref{tab:interventions} treats language dependence and appearance dependence
as separate experimental axes. Swapping the query on an otherwise identical
board asks whether the point follows the requested object. Swapping view positions
asks whether the model can preserve that binding away from a preferred location.
The paired material comparison then asks whether a different appearance changes
the same binding. This structure makes a successful point more informative than
success on a single large, centered target.

Material conversion changes which queries succeed. Qwen3-VL-8B
solves \LargeFullOrig\ original category queries and \LargeFullConv\ converted queries; the image-center
baseline solves none. Qwen3-VL-2B has \SmallFullOrig\ and \SmallFullConv\ hits in the two conditions but changes hit
status on \SmallFullFlips\ paired queries. Gemma4 moves from \GemmaFullOrig\ to \GemmaFullConv\ hits with \GemmaFullFlips\
transitions. As in navigation, equal or nearly equal totals need not identify
the same successful decisions. The action-framed query changes this pattern:
Qwen3-VL-2B moves from \SmallActionOrig\ to \SmallActionConv\ hits, whereas Qwen3-VL-8B retains \LargeActionOrig\
in both conditions. The gain for the smaller model is concentrated in action-framed queries,
showing that the response to appearance depends on how the object is requested.

\paragraph{Scale and model behavior.}
The broader paired archive supplies a complementary readout experiment over six
model configurations, including close and wider views. Its corrected tables
report end-to-end outcomes and output coverage together
(Appendix~\ref{app:archive}). These data expose large differences between a
generated number and a usable spatial prediction. For example, a low hit rate
under one requested frame can coexist with much better spatial alignment under
another interpretation of the same response. The two-view task tests the visual
and language binding that remains after a frame and parser have been fixed.
